\section{Evaluation}
\label{sec:eval}

% SCOPE: the reasoner's own throughput and footprint. Numbers from a fresh
% Apple M4 run of rust/anomaly-addon/examples/edge_bench. NO BEAM, no end-to-end
% pipeline. Humanizer: plain voice, no em/en dashes.

\paragraph{Scope.} We measure the reasoner in isolation: how many sample
evaluations one core sustains, and how much memory the per-series state costs.
End-to-end throughput is set by upstream sample extraction and downstream I/O,
which this paper does not study.

\subsection{Methodology}
The harness is \texttt{rust/anomaly-addon/examples/edge\_bench.rs}. It builds a
\texttt{DetectorEngine} with the default configuration (window 300, minimum 30
samples, 3-sigma threshold, 5-slot confirmation), pre-builds the series keys so
string allocation is not on the measured path, and drives a fixed number of
samples per series with deterministic per-series noise around a baseline of 100
and a rare injected spike once the window is warm. Each call to
\texttt{engine.evaluate} runs the full per-series path: rebuild the rolling
accumulator from the window, evaluate the signal, apply confirmation hysteresis,
and produce a verdict. Decode and transport are not included.

Measurements were taken on an Apple M4 (10 cores), \texttt{rustc} 1.93.1, release
build, single thread, three million evaluations per run. Maximum resident set size
is from \texttt{/usr/bin/time -l}. \NEEDCITE{re-measure throughput and RSS on the
Linux x86 deployment target before submission; macOS over-reports RSS, and the
cgroup budget is enforced on Linux.}

\subsection{Throughput}
\begin{table}[t]
  \centering
  \caption{Reasoner throughput on a single core (Apple M4), three million
  evaluations per run. Throughput tracks how full the window is, not the series
  count: the 50{,}000-series run completes only 60 rounds, so its windows never
  fill and each evaluation rebuilds a shorter accumulator.}
  \label{tab:reasoner-perf}
  \begin{tabular}{rrlrr}
    \toprule
    Series & Rounds & Window state & Throughput (eval/s) & Max RSS \\
    \midrule
    2{,}000  & 1{,}500 & saturated (300) & 479{,}000   & 34.7 MiB \\
    5{,}000  & 600     & mostly full     & 510{,}000   & 81.8 MiB \\
    10{,}000 & 300     & filling         & 646{,}000   & 160 MiB  \\
    50{,}000 & 60      & cold ($<$60)    & 1{,}365{,}000 & 160 MiB \\
    \bottomrule
  \end{tabular}
\end{table}

\Cref{tab:reasoner-perf} reports the sweep. The steady-state figure is the
2{,}000-series run, whose windows stay saturated at 300 samples for most of the
run: a single core sustains about 480{,}000 evaluations per second. Throughput
rises as windows shorten because the shipped path rebuilds the Welford
accumulator from the window on each evaluation, so per-sample cost grows with
window length. An agent emits far fewer than 480{,}000 samples per second, so one
core has ample headroom for real-time detection.

\subsection{The cost of statelessness}
The shipped engine persists only the bounded window tail per series, not a running
accumulator, and rebuilds the accumulator from the window on each evaluation. This
is an $O(\text{window})$ recompute rather than the $O(1)$ incremental update that a
reversible Welford accumulator allows; the shared core implements both. A
microbenchmark that isolates the reasoner from the per-series engine
(\texttt{anomaly-core/examples/reasoner\_bench.rs}, same machine, $300$-sample
window, five million iterations) puts numbers on the tradeoff
(\cref{tab:reasoner-decomp}). The arithmetic kernel itself runs at about $116$
million evaluations per second with the $O(1)$ reversible accumulator, the bare
$O(\text{window})$ recompute at about $0.98$ million, and the full shipped
\texttt{reason\_impl} (the \texttt{CausalFlow} pipeline, a fresh context, and the
allocated verdict) at about $0.46$ million, which matches the full-engine figure of
\cref{tab:reasoner-perf}. The stateless design gives up more than two orders of
magnitude of raw kernel throughput, and the result is still far more than an agent
produces.

\begin{table}[t]
  \centering
  \caption{Reasoner microbenchmark on a single core (Apple M4), $300$-sample
  window, five million iterations, isolating the reasoner from the per-series
  engine.}
  \label{tab:reasoner-decomp}
  \begin{tabular}{lr}
    \toprule
    Path & Throughput (eval/s) \\
    \midrule
    Kernel, incremental $O(1)$ (reversible Welford)      & 116{,}000{,}000 \\
    Kernel, recompute $O(W)$ (rebuild + score)           & 985{,}000 \\
    \texttt{reason\_impl} (\texttt{CausalFlow}, shipped) & 464{,}000 \\
    \bottomrule
  \end{tabular}
\end{table}

The engine uses the stateless path on purpose: the per-series state stays tiny and
trivially serializable for the restart checkpoint, and the reasoner never trusts an
accumulator whose stored summary might disagree with the window it is paired with.

\subsection{Memory footprint}
Per-series state is one window of up to 300 \texttt{f64} values plus a small
counter, so the working-set floor is about 2.6 KiB per series. \cref{tab:reasoner-perf}
shows higher resident sizes on macOS because the system allocator over-reports.
\NEEDCITE{the deployment target is Linux under a 256 MiB / 50\% CPU cgroup budget;
re-measure RSS there.} \TODO{at 50{,}000 series with full windows, memory rather
than CPU is the binding constraint; report where the series cap and the cgroup
budget meet.}

\subsection{Threats to validity}
\dossier{no detection-accuracy (precision/recall) evaluation; synthetic workload,
not production traces; single machine, single run, no variance reported; macOS is
not the deployment platform; $z$-score assumes approximate local stationarity;
point anomalies only at this timescale.} \TODO{}
